\subsection{Maximum Likelihood Detector}\label{subsubsec:ml}
% FONTE: TG1 (sbc-template.tex), Equalization.
The optimal method for recovering transmitted symbols in MIMO systems is the
Maximum Likelihood (ML) detector \cite{wang2021tlbdinsa}. The algorithm estimates the transmitted
symbols by finding
the candidate that minimizes the Euclidean distance between the received
signal $\mathbf{y}$ and $\mathbf{H}\mathbf{s}$:

\begin{equation}
\hat{\mathbf{s}} = \underset{\mathbf{s} \in \mathcal{S}^{K}}{\arg\min}
\left\|\mathbf{y} - \mathbf{H}\mathbf{s}\right\|_2^2
\label{eq:ml}
\end{equation}
where $\mathcal{S}^{K}$ is the set of all possible combinations of
constellation symbols transmitted by the $K$ users, and
$\left\|\cdot\right\|_2^2$ denotes the sum
of the squared magnitudes of all entries of a vector (squared Euclidean norm).

\textcolor{blue}{This method yields the lowest possible bit error rate without channel coding.} However, MIMO detection is NP-hard \cite{zero-forcing}, and an exhaustive search over $\mathcal{S}^{K}$ has a complexity of $\mathcal{O}(|\mathcal{S}|^K)$, making the ML detector impractical for systems with a large number of
users or high-order modulation schemes~\cite{bello2022algorithms}.
% FONTE: wang2021tlbdinsa p. 2971.

\subsection{Filters}\label{subsec:filtros}\label{subsubsec:wiener}

A filter receives a measured signal and produces another signal at its output. In a linear filter, each output sample is a linear combination of the measurements, where the weights serve as the filter coefficients \cite{oppenheim2010wiener}.

Linear equalizers are linear filters and suboptimal alternatives to the ML detector whose
complexity is independent of the constellation scheme employed. These
techniques estimate a continuous-valued approximation $\tilde{\mathbf{s}}$
(soft-output) of $\mathbf{s}$ by applying a weight matrix $\mathbf{W}$ to
the received signal vector $\mathbf{y}$~\cite{rodriguezsanchez2020}:
\begin{equation}\label{eq:filtro}
\tilde{\mathbf s}=\mathbf W\mathbf y.
\end{equation}

The $k$-th row of $\mathbf W$ contains the weights used to estimate the symbol for UE $k$. Linear equalizers differ in the criterion used to select these weights \cite{tse-viswanath}.
% FONTE: oppenheim2010wiener p. 195, Eq. 11.1; tse-viswanath pp. 351 e 358.

\subsubsection{Zero-Forcing Equalizer}\label{subsubsec:zf}
The Zero-Forcing (ZF) equalizer aims to eliminate interference between the UEs, without considering noise \cite{tse-viswanath}. To do this, it inverts the channel: multiplying $\mathbf y$ by the inverse of $\mathbf H$ undoes the mixing of the UE signals and leaves only the corresponding UE symbol at each output. However, $\mathbf H$ is an $M \times K$ matrix, which is not square and has no inverse. \textcolor{blue}{ZF uses the pseudo-inverse of $\mathbf H$ instead, which has the same effect of separating the symbols \cite{zero-forcing}:}
% FONTE: tse-viswanath pp. 351-352 e 358.
\begin{equation}
\mathbf{W}_\text{ZF}=
\left(\mathbf{H}^H\mathbf{H}\right)^{-1}\mathbf{H}^H,
\label{eq:zf_filter}
\end{equation}
where $\mathbf{H}^H$ denotes the conjugate transpose of $\mathbf{H}$. The product $\mathbf H^H\mathbf H$ is the Gram matrix of $\mathbf H$ \cite{jeon2019}. It is square, of size $K \times K$, and it is this matrix that ZF inverts. \textcolor{blue}{Noise also passes through the filter, since it is not considered.} As a consequence, ZF amplifies the noise present at the receiver, particularly in poorly conditioned channels.
% FONTE: jeon2019 L105 (Gram matrix).

\subsubsection{LMMSE Equalizer}\label{subsubsec:lmmse}
The Linear Minimum Mean Square Error (LMMSE) equalizer addresses the ZF
limitation by explicitly accounting for the noise statistics during equalizer
design \cite{tse-viswanath}. The estimate $\tilde{\mathbf{s}}_\text{LMMSE}$ is obtained by applying the
equalization matrix to the received signal:

\begin{equation}
\tilde{\mathbf{s}}_\text{LMMSE}=
\mathbf{W}_\text{LMMSE}\mathbf{y}
\label{eq:lmmse_estimate}
\end{equation}

The weights of $\mathbf W$ are chosen according to the criterion:

\begin{equation}
\min_{\mathbf{W}}
\mathbb{E}
\left[
\left\|\mathbf{W}\mathbf{y}-\mathbf{s}\right\|_2^2
\right],
\label{eq:lmmse_opt}
\end{equation}
% FONTE: zhao2025jsac paper.md:67 (Eq. 2); haykin p. 198, Eq. 5.36.
which minimizes the mean squared error of the equalization \cite{zhao2025jsac}. This error is calculated using the expected value $\mathbb E[\cdot]$, the average value of a random variable, since the symbols and noise are random \cite{haykin}. Thus, the minimization considers the different realizations of these variables, and not just a specific observation.
The squared error includes both interference and noise. \textcolor{blue}{Therefore, LMMSE can retain some interference to reduce noise, unlike ZF \cite{tse-viswanath}.}

% FONTE: tse-viswanath p. 358.
Assuming that the channel matrix $\mathbf{H}$ is perfectly known, the
corresponding solution for the LMMSE equalizer in a centralized MIMO system
is given by~\cite{zhao2025jsac}:

\begin{equation}
\mathbf{W}_\text{LMMSE}=
\left(
\mathbf{H}^H\mathbf{R}_n^{-1}\mathbf{H}
+
\frac{1}{E_s}\mathbf{I}
\right)^{-1}
\mathbf{H}^H\mathbf{R}_n^{-1}
\label{eq:lmmse_solution}
\end{equation}
where $\mathbf{n}$ is modeled as a zero-mean random vector with covariance
matrix $\mathbf{R}_n = \mathbb{E}[\mathbf{n}\mathbf{n}^H]$ (Section~\ref{subsubsec:interferencia}), and $E_s$
denotes the average energy per transmitted symbol. \textcolor{blue}{The LMMSE approaches the ZF with low noise \cite{tse-viswanath}.}
% FONTE: tse-viswanath p. 361.

The LMMSE equalizer leads to a slight increase in computational complexity
compared with ZF, as it additionally requires estimating the noise covariance matrix $\mathbf{R}_n$ to compute the filter. Both equalizers share the same asymptotic complexity
$\mathcal{O}(K^3)$, dominated by the matrix inversion, but the LMMSE
generally achieves a lower BER than ZF~\cite{li-tradeofs}.
% FONTE: li-tradeofs Fig. 5(a).

The equalization
matrix $\mathbf{W}$ is computed only once per coherence block and reused for
each received symbol, reducing the average computational cost per symbol by a
factor of $1/N_\text{coh}$, where $N_\text{coh}$ is the number of symbols in the block \cite{zhao2025jsac}.
% FONTE: zhao2025jsac paper.md:561.
